\documentclass[twocolumn]{book}
\usepackage[a4paper,margin=2cm]{geometry}
\usepackage{background}
\usepackage{titlesec}
\usepackage{xcolor}
\usepackage{titling}
\usepackage{tocloft} % optional: control TOC formatting
\usepackage{parskip}
\usepackage{setspace}
\usepackage{graphicx}
\usepackage{xspace}
% Optional: Use a special font for the title
\usepackage{tgchorus} % Elegant script-style font
\usepackage{cmbright}
\usepackage{hyperref}
\hypersetup{pdfborder=0 0 0}
\usepackage{xparse}
\usepackage{calc}
\usepackage{shadowtext}
\usepackage{tcolorbox}
\usepackage{indentfirst}
\usepackage{tikz}
\usetikzlibrary{shadows.blur}
\usepackage[utf8]{inputenc}


% Column spacing
\setlength{\columnsep}{20pt}

% Define dark blue color
\definecolor{darkblue}{RGB}{0, 51, 102}

% Section styling
\titleformat{\section}
  {\bfseries\Large\color{darkblue}}
  {\thesection}{1em}{}
\titlespacing*{\section}{0pt}{36pt}{24pt}

\titleformat{\subsection}
  {\bfseries\large\color{darkblue}}
  {\thesubsection}{1em}{}
\titlespacing*{\subsection}{0pt}{24pt}{18pt}

% Line spacing
% \renewcommand{\baselinestretch}{1.5}

\definecolor{marineblue}{RGB}{0, 51, 102}


\newcommand{\setTitleBackground}[2]{%
  \backgroundsetup{
    scale=1,
    angle=0,
    opacity=#2,
    contents={\includegraphics[width=\paperwidth,height=\paperheight]{#1}}
  }%
}

\usepackage{fancyhdr}
\pagestyle{fancy}

% Stocker le titre du chapitre
\newcommand{\chaptertitle}{}

% Style pour la 2e page du chapitre
\fancypagestyle{chapterstyle}{
  \fancyhf{}
  \fancyhead[L]{\thechapter.\ \chaptertitle} % numéro + titre
  \fancyhead[R]{\thepage}                    % numéro de page
  \renewcommand{\headrulewidth}{0pt}
}

% Style normal pour le reste du document
\fancypagestyle{normalstyle}{
  \fancyhf{}
  \fancyhead[R]{\thepage}
  \renewcommand{\headrulewidth}{0pt}
}


\NewDocumentCommand{\NgelmakProjectNewChapter}{m o o}{%
  % Images optionnelles
  \IfNoValueTF{#2}{\def\second{images/marios_gkortsilas_jceqhptpu24_unsplash.jpg}}{\def\second{#2}}%
  \IfNoValueTF{#3}{\def\third{images/georgie_cobbs_muohbrfgeqy_unsplash.jpg}}{\def\third{#3}}%

  % Passage en une colonne
  \clearpage
  \onecolumn

  % Numérotation + TOC
  \refstepcounter{chapter}
  \addcontentsline{toc}{chapter}{#1}

  % Stocker le titre pour l'en-tête
  \renewcommand{\chaptertitle}{#1}

  % --- PAGE 1 DU CHAPITRE ---
  \thispagestyle{empty}
  \setTitleBackground{\second}{1.0}

  \vspace*{\fill}
  \begin{center}
    \begin{tikzpicture}
      \node[
        draw=none,
        fill=white,
        fill opacity=0.5,
        text opacity=1,
        anchor=west,
        inner sep=2cm,
        minimum width=\textwidth
      ] at (0,0) {
        \begin{minipage}{\textwidth}
          \raggedright
          \setstretch{1.4}
          {\fontsize{28}{34}\selectfont\bfseries\color{black} #1}
        \end{minipage}
      };
    \end{tikzpicture}
  \end{center}
  \vspace*{\fill}

  % --- PAGE 2 DU CHAPITRE ---
  \clearpage
  \thispagestyle{chapterstyle}
  \setTitleBackground{\third}{0.1}

  % Retour au mode deux colonnes
  \clearpage
  \twocolumn

  % --- IMPORTANT : appliquer le style de chapitre à toutes les pages suivantes ---
  \pagestyle{chapterstyle}
}


\newcommand{\ngelmak}{\textsc{Ngelmak}\xspace}



\begin{document}
\pagestyle{normalstyle}

\setTitleBackground{images/kasra_askari_5v9rtr8qp8_unsplash.jpg}{1.0}

\vspace*{3cm}

% Title — centered across full width
\noindent
\begin{minipage}{\textwidth}
  \centering
  \begin{tikzpicture}
    \node[
      rounded corners=3mm,
      fill=black,
      fill opacity=0.8,
      text opacity=1,
      text=white,
      align=center,
      anchor=center,   % <<< THIS keeps the box centered
      inner xsep=40pt, % horizontal padding
      inner ysep=15pt, % vertical padding
    ] {
      \fontsize{48pt}{64pt}\selectfont\bfseries Projet Ngelmak
    };
  \end{tikzpicture}
\end{minipage}

\vspace{0.5cm}

% Subtitle — new line, aligned right
\begin{minipage}{\textwidth}
  \raggedleft
  \begin{tikzpicture}
    \node[
      rounded corners=6pt,
      fill=black,
      fill opacity=0.8,
      text opacity=1,
      text=white,            % <<< CHANGE TEXT COLOR HERE
      inner sep=5pt
    ] {\large\textbf{Soutenir l'autonomie grâce au partage et à l'entraide}};
  \end{tikzpicture}
\end{minipage}

\vspace*{\fill}

\noindent
\makebox[\textwidth]{%
	\begin{minipage}[b]{0.6\textwidth}
		\raggedright
		{\large\color{white} Youssouph Faye \quad} \\
		{\normalsize\color{white} October 16, 2025}
	\end{minipage}%
	\hfill
	\begin{minipage}[b]{0.4\textwidth}
		\raggedleft
		\includegraphics[width=2.5cm]{images/Ngelmak-Project/ngelmak-logo-light.pdf}
	\end{minipage}%
}

\clearpage
\newpage

% \vspace*{3cm}
\onecolumn
\setTitleBackground{images/georgie_cobbs_muohbrfgeqy_unsplash.jpg}{0.1}
\chapter*{Résumé}
\Large L'Afrique évolue dans un paysage informationnel dominé par des plateformes étrangères qui imposent leurs récits, filtrent nos voix et façonnent la perception de nos réalités, au point que la désinformation et la propagande anti-africaine sont devenues quotidiennes. Cette dépendance, aggravée par l'absence d'alternatives souveraines proposées par nos dirigeants ou nos institutions, prive les Africains d'un espace où communiquer librement, documenter leurs expériences et construire leur propre narratif. Face à ce vide stratégique, la création d'un réseau africain autonome apparaît comme une nécessité vitale pour reprendre la parole, protéger notre mémoire collective et affirmer notre souveraineté. Le projet \ngelmak s'inscrit dans cette dynamique : bâtir un outil modeste au départ, mais évolutif, crédible et durable, qui permette enfin aux Africains de s'exprimer, d'analyser et de débattre sans dépendre de structures hostiles à leurs intérêts.


\clearpage
\backgroundsetup{contents={},opacity=0} % no content, invisible
\tableofcontents
\clearpage


% ** A NEW CHAPTER **
\NgelmakProjectNewChapter{Vision Générale, Motivation et Fondements du Projet \ngelmak}

\section{Contexte géopolitique et enjeux de la communication africaine}

La communication occupe aujourd'hui une place centrale dans les dynamiques géopolitiques mondiales. Pourtant, nous, Africains, continuons de dépendre de canaux d'information conçus et contrôlés par des acteurs extérieurs, qui imposent leur propre narratif sur notre continent. Cette dépendance structurelle leur permet de façonner la perception des événements qui se déroulent chez nous, au point que nous nous retrouvons souvent noyés dans la désinformation, la manipulation et la propagande anti-africaine. Ce phénomène est devenu si courant qu'il ne nécessite même plus d'exemples : il constitue désormais une réalité quotidienne pour des millions d'Africains qui, parfois, connaissent à peine leurs propres terres et leurs propres histoires.

À cette domination narrative s'ajoute une forme de censure systématique. Dès que nous proposons un récit différent de celui imposé par des acteurs hostiles à nos intérêts, nos voix sont filtrées, marginalisées ou réduites au silence. Ce déséquilibre est aggravé par l'absence d'alternatives crédibles proposées par nos dirigeants, nos institutions régionales ou continentales. Ni les États, ni les organisations sous-régionales, ni même l'Union africaine ne mettent en place des solutions souveraines permettant aux Africains de communiquer entre eux, avec leurs peuples ou avec le reste du monde. Pourtant, de nombreuses technologies open source existent déjà et pourraient être adaptées à nos besoins.

Malgré cela, nous continuons à utiliser des plateformes étrangères pour exprimer nos opinions, partager nos idées ou débattre de nos enjeux. Cette dépendance nous fragilise, car elle place notre parole, notre mémoire collective et notre capacité d'analyse entre les mains d'acteurs dont les intérêts ne coincident pas nécessairement avec les nôtres. Dans un monde où l'information est devenue un instrument de pouvoir, cette situation constitue un handicap majeur pour notre souveraineté.

\section{Nécessité d'un réseau souverain et motivations du projet \ngelmak}

Face à ce constat, la nécessité de construire un réseau conçu par des Africains et destiné aux Africains apparaît avec évidence. C'est dans cet esprit qu'est né le Projet \ngelmak. \ngelmak, qui signifie \textit{cour} en sérère, renvoie à cet espace traditionnel de rassemblement où l'on échange, débat et construit ensemble, et c'est précisément ce rôle que nous voulons redonner à notre espace numérique souverain. Nous devons disposer d'un environnement capable de transmettre notre narratif, de documenter notre présent et de préserver notre mémoire collective. Un espace où nous pouvons partager nos idées, analyser nos réalités et débattre de nos enjeux sans craindre la censure, la condescendance ou la manipulation.

Ce projet vise à nous sortir de la posture réactive dans laquelle nous sommes enfermés depuis trop longtemps. Nous passons notre temps à dénoncer le racisme, la manipulation médiatique, l'ingérence ou la désinformation, sans jamais disposer d'un outil qui nous permette de proposer notre propre vision du monde. Les technologies existent déjà ; il nous appartient désormais de les utiliser pour créer, à travers \ngelmak, un environnement numérique qui serve nos intérêts et reflète nos valeurs.

Nous sommes conscients que le démarrage sera modeste et vulnérable. Nous ne rivaliserons pas immédiatement avec les grandes plateformes mondiales. Nous affronterons des difficultés, nous connaîtrons des échecs, mais nous apprendrons, nous progresserons et nous nous renforcerons. Ce processus est nécessaire pour bâtir un outil durable, crédible et utile à long terme. Notre détermination doit être totale, car ce projet s'inscrit dans une démarche de souveraineté, de dignité et de responsabilité collective.

Cette initiative dépasse nos ressources individuelles, mais elle s'inscrit dans une volonté commune de contribuer à la renaissance africaine. Nous devons cesser de chercher la reconnaissance dans des institutions qui ne nous servent pas. Nous devons nous imposer par des réalisations concrètes, par des faits, par des outils comme \ngelmak, qui démontrent notre capacité à construire, à innover et à nous organiser. L'Afrique dispose de tout le potentiel nécessaire pour avancer ; il nous revient de l'activer.

\section{Vision panafricaine, unité et principes éthiques}

\ngelmak s'inscrit dans une vision panafricaine fondée sur l'unité, la dignité et la souveraineté intellectuelle. Nous sommes conscients que la majorité des pays africains ont été artificiellement créés par des puissances étrangères, sans lien avec nos peuples ni nos réalités. Cette fragmentation imposée a laissé des traces profondes, tant physiques que spirituelles. Nous refusons que cette division historique continue de nous séparer dans notre manière de penser, de communiquer ou de nous percevoir collectivement.

Malgré nos différences idéologiques, linguistiques, religieuses, culturelles ou sociales, nous partageons une condition commune : nous sommes souvent perçus de l'extérieur comme un ensemble homogène, réduit à des stéréotypes déshumanisants. Ce projet n'a pas pour objectif de convaincre qui que ce soit du contraire, mais de nous permettre de nous rassembler, de construire notre propre espace et de produire notre propre récit. Il s'agit de reprendre la maîtrise de notre parole, de notre image et de notre histoire.

Dans cette perspective, nous établissons des principes éthiques clairs. Aucune forme de condescendance, d'insulte ou de dénigrement ne sera tolérée. Les expressions visant à réduire au silence, à diviser ou à humilier — telles que "occupez-vous de votre pays" ou toute variante similaire — n'auront pas leur place dans cet espace. Nous voulons créer un environnement où chacun peut s'exprimer avec respect, rigueur et responsabilité, dans l'intérêt collectif.

Nous affirmons également que notre projet doit permettre aux Africains de parler de l'Afrique, mais aussi du monde entier. Les dynamiques globales nous concernent, car elles influencent directement nos sociétés, nos économies et nos trajectoires. Nous ne devons pas nous limiter à commenter ce qui se passe chez nous : nous devons participer pleinement aux discussions internationales, avec nos propres analyses et nos propres perspectives.

Cette vision éthique et panafricaine constitue le socle sur lequel repose l'ensemble du projet \ngelmak. Elle guide nos choix, nos priorités et notre manière de concevoir les outils que nous mettons en place. Elle nous rappelle que notre objectif n'est pas seulement technique ou fonctionnel, mais profondément humain, culturel et politique au sens noble du terme : contribuer à l'éveil, à la cohésion et à la souveraineté de nos peuples.

\section{Autonomie et financement du projet \ngelmak}

Le financement du projet \ngelmak repose sur une approche simple, transparente et entièrement indépendante. Le développement initial est assuré par les ressources personnelles du fondateur et des membres les plus investis. Ce choix traduit la volonté de préserver l'autonomie du projet et d'éviter toute influence extérieure susceptible d'en altérer les valeurs ou la vision.

À mesure que la plateforme évoluera, des dons volontaires pourront être sollicités auprès de la communauté. Ces contributions resteront toujours facultatives et ne conditionneront jamais l'accès aux services. Elles auront pour seule finalité de couvrir les coûts techniques, matériels et humains nécessaires au maintien et à l'amélioration de l'outil. Aucun service payant, aucune fonctionnalité premium et aucune forme de commission ne seront proposés. Le projet n'a pas vocation à générer un profit : il repose sur un modèle de solidarité et de participation collective.

Cette approche garantit que l'ensemble des services restera accessible à tous, sans distinction de moyens, et que le projet demeurera libre de toute pression économique, politique ou commerciale. L'objectif n'est pas de créer une entreprise, mais de bâtir un outil communautaire, durable et au service de l'intérêt général.



\section{Fondements sacrés du projet \ngelmak}

\ngelmak repose sur un ensemble de principes intangibles, considérés comme sacrés et non négociables. Ces fondements définissent son identité, orientent sa gouvernance et garantissent sa fidélité à la vision qui l'a fait naître. Ils constituent la base éthique et politique sur laquelle repose l'ensemble de l'initiative.

\paragraph{Gratuité absolue.}
Tous les services proposés par le projet resteront toujours gratuits, sans aucune exception. La formule souvent répétée — "si c'est gratuit, c'est vous le produit" — ne s'applique pas ici. Nous rejetons toute forme de monétisation directe ou indirecte des utilisateurs : aucune fonctionnalité ne sera payante, aucune commission ne sera prélevée, et aucun accès ne sera conditionné à un abonnement ou à un paiement. La gratuité n'est pas un argument commercial, mais un engagement éthique fondamental, garantissant que l'outil demeure un bien commun accessible à tous.

\paragraph{Absence totale de volonté de profit.}
\ngelmak n'a pas vocation à générer des revenus, à attirer des investisseurs ou à se transformer en entreprise lucrative. Il s'agit d'un outil communautaire conçu pour répondre à des besoins collectifs, et non pour servir des intérêts privés. Cette orientation exclut toute forme de publicité intrusive, de vente de données, de partenariats commerciaux contraires à nos valeurs ou de modèle économique visant à tirer profit des utilisateurs. \ngelmak ne poursuivra jamais un objectif financier : il est et restera non lucratif.

\paragraph{Protection stricte des données.}
Aucune donnée ne sera vendue, partagée ou exploitée par des tiers, sous quelque forme que ce soit. Les informations confiées par les utilisateurs ne seront ni échangées, ni utilisées à des fins commerciales, politiques ou publicitaires. La confidentialité et la sécurité des données constituent des engagements absolus, respectés à chaque étape du développement, du fonctionnement et de l'évolution du projet. La protection des utilisateurs n'est pas une fonctionnalité : c'est un principe fondateur.

\paragraph{Gouvernance centrée sur la communauté.}
Toute évolution majeure des politiques, des règles fondamentales ou des orientations générales devra être soumise à la communauté. \ngelmak n'appartient pas seulement à ses fondateurs : il appartient à celles et ceux qui l'utilisent et le font vivre. Cette gouvernance participative garantit que les décisions importantes ne pourront jamais être prises de manière unilatérale ou contre l'intérêt collectif. Elle assure la continuité du projet, sa stabilité et sa fidélité à ses valeurs initiales.

\paragraph{Transparence totale.}
Les décisions importantes, les évolutions techniques, les orientations stratégiques et les choix de gouvernance seront documentés et accessibles à tous. Cette transparence s'applique également à l'usage des dons : chaque contribution financière fera l'objet d'un suivi clair, public et vérifiable, permettant à la communauté de connaître précisément l'affectation des ressources. Rien de ce qui concerne le fonctionnement, le financement ou l'avenir du projet ne sera dissimulé ou décidé dans l'ombre. La transparence constitue un pilier essentiel de la confiance, de la légitimité et de la pérennité du projet.


Ces fondements sacrés définissent ce que nous sommes, ce que nous voulons construire et ce que nous refusons de devenir. Ils garantissent que \ngelmak restera un espace libre, éthique, souverain et profondément ancré dans l'intérêt général.


\section{Droits, propriété et gouvernance du projet}

\ngelmak repose sur un cadre juridique et organisationnel destiné à garantir son ouverture, sa pérennité et la protection de sa vision fondatrice. L'ensemble du code source, de la documentation, de l'architecture et des éléments conceptuels a été initié, conçu et développé par le fondateur du projet, qui en détient les droits d'auteur initiaux. Ces droits constituent la base légale du projet et définissent son identité.

Le projet \ngelmak est publié sous une licence open source de type GNU, garantissant que le code restera libre, accessible et modifiable par tous. Cette licence protège \ngelmak contre toute tentative de privatisation, d'appropriation ou de détournement. Elle assure également que toute contribution future restera compatible avec les principes de liberté, de transparence et de partage qui fondent le projet.

La gouvernance du projet repose sur trois niveaux complémentaires.

\paragraph{Fondateur.}
Le fondateur est le détenteur du copyright initial et le garant de la vision. Il conserve l'autorité finale sur les orientations stratégiques, la validation des versions officielles, la définition des règles fondamentales et la protection de l'identité du projet. Ce rôle est permanent et ne peut être retiré, transféré ou remis en cause par aucune instance interne ou externe. Le fondateur demeure le gardien ultime de la cohérence, de la continuité et de la fidélité du projet à ses fondements sacrés.

\paragraph{Contributeurs.}
Les contributeurs regroupent toutes les personnes participant au développement, à la rédaction, à la conception ou au financement. Leur participation est reconnue, valorisée et intégrée dans le projet. Ils peuvent proposer des améliorations, développer des modules, rédiger des documents ou contribuer à la maintenance. Toutefois, leurs contributions s'inscrivent dans le cadre défini par les principes fondateurs, et toute modification majeure nécessite la validation du fondateur.

\paragraph{Communauté.}
La communauté occupe une place centrale dans la gouvernance. Elle peut proposer des idées, voter sur certaines orientations, signaler des problèmes, participer à la réflexion collective et contribuer à l'évolution du projet. Aucun changement de politique, aucune modification des fondements sacrés et aucune évolution majeure ne pourra être adoptée sans consultation de la communauté. Toutefois, cette participation ne remet pas en cause le rôle du fondateur, qui demeure l'autorité ultime chargée de préserver la cohérence et l'intégrité du projet.

Ce cadre de propriété et de gouvernance assure que \ngelmak restera un bien commun, libre, éthique et durable, tout en protégeant la vision initiale contre toute tentative d'appropriation ou de déviation.


\section{Clarification des décisions stratégiques}

Dans le cadre de ce projet, une décision stratégique désigne toute décision susceptible d'affecter la vision, l'identité ou les fondements du projet. Cela inclut notamment :

\begin{itemize}
    \item la modification des principes sacrés (gratuité, absence de profit, protection des données) ;
    \item l'introduction de publicité, de commissions ou de services payants ;
    \item tout changement de licence ou de modèle économique ;
    \item toute tentative de privatisation, de commercialisation ou de mise en bourse ;
    \item toute modification de la gouvernance ou des droits fondamentaux ;
    \item toute évolution susceptible de détourner \ngelmak de ses objectifs initiaux.
\end{itemize}

Ces décisions nécessitent obligatoirement la consultation de la communauté, la validation du fondateur et le respect strict des fondements sacrés. Le rôle du fondateur n'est pas un pouvoir absolu, mais un pouvoir de protection destiné à garantir que \ngelmak ne pourra jamais être détourné, privatisé ou transformé en produit commercial. Les décisions techniques, opérationnelles ou fonctionnelles restent ouvertes à la participation collective.


% \section{Perspectives d'évolution de la gouvernance}

% Pour renforcer la structure du projet et garantir sa pérennité, plusieurs éléments pourront être développés dans les phases ultérieures :

% \begin{itemize}
%     \item une charte de gouvernance complète définissant les rôles et responsabilités ;
%     \item une charte éthique détaillant les valeurs et les règles de comportement ;
%     \item un document de contribution (CONTRIBUTING) décrivant les processus de participation ;
%     \item un statut de projet communautaire formalisant les droits et devoirs de chacun ;
%     \item une licence GNU adaptée, éventuellement complétée par des clauses éthiques supplémentaires.
% \end{itemize}

% Ces outils permettront d'assurer la continuité, la transparence et la cohérence du projet tout en renforçant la participation de la communauté.


\section{Conclusion}

Ce chapitre a posé les bases conceptuelles, éthiques et organisationnelles du projet {\ngelmak{}}. Il a présenté les motivations profondes qui justifient sa création, les enjeux géopolitiques qui rendent indispensable l'émergence d'un espace numérique souverain, ainsi que les principes sacrés qui guideront son développement. \ngelmak s'inscrit dans une démarche de responsabilité collective, de transparence et de souveraineté, en affirmant clairement son refus de toute forme de privatisation, de commercialisation ou de compromission avec des intérêts extérieurs.

La structure de gouvernance définie ici garantit un équilibre entre l'ouverture à la participation communautaire et la protection de la vision fondatrice. Le rôle du fondateur, les responsabilités des contributeurs et la place centrale de la communauté ont été établis de manière à assurer la continuité, la cohérence et la fidélité du projet à ses valeurs. Les mécanismes de décision, en particulier ceux concernant les orientations stratégiques, ont été conçus pour empêcher toute dérive susceptible de détourner \ngelmak de ses objectifs initiaux.

Les fondements sacrés — gratuité absolue, absence de profit, protection stricte des données, transparence et gouvernance communautaire — constituent le socle immuable sur lequel reposera l'ensemble des outils et services développés. Ils garantissent que \ngelmak demeurera un bien commun numérique, accessible à tous et protégé contre toute tentative de déviation ou d'appropriation.

Ce cadre général ouvre la voie aux développements techniques, organisationnels et communautaires qui seront détaillés dans les chapitres suivants. Il marque le point de départ d'une initiative ambitieuse, construite avec rigueur et conviction, et destinée à servir durablement les intérêts des Africains et de tous ceux qui partagent cette vision.

% ** A NEW CHAPTER **
% \NgelmakProjectNewChapter{Un Espace d'Expression Responsable et Constructive}


% \section{Introduction}

% Le chapitre précédent a posé les fondements éthiques, philosophiques et organisationnels de notre projet. Nous avons affirmé notre volonté de créer un espace souverain, responsable et respectueux, où la parole des citoyens peut s'exprimer librement sans être déformée, censurée ou instrumentalisée. Le partage de l'information constitue le cœur de cette ambition. Il ne s'agit pas de reproduire les dynamiques superficielles, individualistes ou sensationnalistes que l'on observe ailleurs, mais de proposer un modèle fondé sur la réflexion, la responsabilité et l'intérêt collectif. Ce chapitre présente la manière dont l'information sera partagée au sein de notre plateforme, ainsi que les principes qui guideront la création, la diffusion et la réception des contenus.

% \section{1. Une philosophie du partage fondée sur l'intérêt public}

% Notre plateforme n'a pas vocation à devenir un espace de mise en scène personnelle, d'auto-promotion ou de recherche de visibilité. Les contenus centrés sur la vie privée, les annonces sans intérêt public, les récits intimes, les provocations, les polémiques artificielles ou les tentatives de sabotage n'y ont pas leur place. Nous refusons les dynamiques où l'on partage pour attirer l'attention plutôt que pour transmettre une idée, une analyse ou une information utile.

% Ici, le partage d'information repose sur une philosophie claire : chaque publication doit apporter quelque chose à la communauté. Elle peut éclairer un sujet, dénoncer une injustice, proposer une réflexion, transmettre un savoir, raconter une expérience significative ou encourager un débat constructif. L'objectif n'est pas de se montrer, mais de contribuer. L'information n'est pas un produit de consommation rapide, mais un outil d'éveil, de compréhension et d'action.

% \section{2. Un cadre d'expression responsable et respectueux}

% Pour garantir la qualité des échanges, nous imposons un cadre strict fondé sur le respect, la rigueur et la responsabilité. Aucune insulte, aucun dénigrement, aucune attaque personnelle ne sera tolérée, qu'elle vise un auteur, un commentateur ou un autre utilisateur. La critique est encouragée, mais elle doit être argumentée, contextualisée et orientée vers la compréhension ou l'amélioration. La plateforme n'est pas un lieu de confrontation, mais un espace de réflexion collective.

% Cette exigence s'applique également aux commentaires. Chaque commentaire doit suivre la même philosophie que les publications : apporter une idée, une nuance, une précision, une critique constructive ou un complément d'information. Les réactions impulsives, les attaques gratuites, les insinuations malveillantes ou les tentatives de discrédit ne seront pas acceptées. Nous voulons créer un environnement où chacun peut s'exprimer sans crainte d'être humilié, agressé ou ridiculisé.

% \section{3. Un système de publication conçu pour la qualité}

% La création d'un compte sera volontairement simple, afin de permettre à chacun de rejoindre la plateforme sans obstacle inutile. Un nom d'utilisateur, un mot de passe et, si l'utilisateur le souhaite, une adresse email suffiront. Une vérification plus avancée sera introduite ultérieurement pour limiter la prolifération de comptes multiples et garantir une utilisation responsable des services. Cette vérification ne sera jamais intrusive et respectera strictement la confidentialité des données.

% Les publications devront respecter un minimum de structure et de clarté. Elles pourront prendre la forme de textes, d'enregistrements vocaux, de vidéos ou de contenus hybrides, mais devront toujours être accompagnées d'une réflexion ou d'un contexte. La plateforme ne valorise pas la quantité, mais la pertinence. Les mécanismes de validation, de modération et de visibilité seront progressivement renforcés grâce à l'implication de volontaires, dont le rôle évoluera au rythme du développement du site.

% \section{4. Une alternative aux modèles existants}

% Notre approche se distingue volontairement des plateformes qui encouragent l'auto-promotion, la mise en scène de soi, la course aux réactions ou la diffusion de contenus sans intérêt public. Nous refusons les logiques de buzz, de compétition sociale ou de voyeurisme. Nous ne cherchons pas à divertir, mais à informer, à éduquer, à éveiller et à soutenir.

% Ici, l'information n'est pas un prétexte pour attirer l'attention, mais un moyen de construire une conscience collective. Les utilisateurs ne sont pas des spectateurs, mais des acteurs. Les publications ne sont pas des vitrines personnelles, mais des contributions à un effort commun. La plateforme n'est pas un réseau social au sens traditionnel, mais un espace d'expression citoyenne, intellectuelle et communautaire.

% \section{Conclusion}

% Le partage de l'information est au centre de notre projet. Il repose sur une vision exigeante, fondée sur l'intérêt public, la responsabilité, le respect et la qualité. Nous voulons offrir un espace où chacun peut s'exprimer librement, mais dans un cadre qui valorise la réflexion plutôt que le bruit, la connaissance plutôt que la distraction, la critique constructive plutôt que l'attaque personnelle. Ce chapitre établit les principes qui guideront la création et la diffusion des contenus, et prépare le terrain pour les mécanismes techniques et organisationnels qui seront détaillés dans les sections suivantes. Notre ambition est de bâtir un outil qui élève, rassemble et éclaire, au service de tous.







\NgelmakProjectNewChapter{Un Outil pour Construire Notre Propre Narratif}

\section{Introduction}

Dans le chapitre précédent, nous avons défini les fondements éthiques et organisationnels du projet, en affirmant notre volonté de bâtir un espace souverain, responsable et orienté vers l'intérêt collectif. Le présent chapitre décrit le cadre dans lequel les membres de la communauté pourront s'exprimer, partager leurs idées et interagir entre eux. Notre objectif n'est pas de reproduire les dynamiques superficielles, individualistes ou sensationnalistes que l'on observe ailleurs, mais de proposer un modèle d'expression fondé sur la réflexion, la responsabilité et la contribution au bien commun. L'expression n'est pas ici un acte de visibilité personnelle, mais un engagement citoyen.

\section{Les Postes et les Interactions Communautaires}


\subsection{Le Poste : unité centrale d'expression}

Le \textit{Poste} constitue l'unité fondamentale d'expression au sein de la plateforme. Il s'agit d'un contenu créé par un utilisateur dans l'intention de partager une idée, une analyse, une critique, un témoignage ou un enseignement. Un Poste peut également répondre à un autre, permettant la construction d'un dialogue structuré et continu entre les membres de la communauté. Par essence, il vise à éclairer, informer ou susciter une réflexion collective, et s'inscrit dans une démarche d'expression citoyenne orientée vers la compréhension et la participation au débat public.

\paragraph{Nature et forme du Poste.}

Un Poste doit présenter une intention claire, une structure minimale et une valeur ajoutée pour la communauté. Il peut prendre différentes formes — texte, enregistrement vocal, vidéo ou contenu hybride — mais doit toujours être accompagné d'un contexte ou d'une réflexion permettant d'en comprendre la portée. Le Poste n'est pas un espace d'auto-promotion, de mise en scène personnelle ou de recherche de visibilité.

\paragraph{Exigences relatives au contenu.}

Pour préserver la qualité de l'espace d'expression, le contenu d'un Poste doit être rédigé avec sérieux, honnêteté et dans un esprit constructif. Il doit être argumenté, factuel lorsque nécessaire, et orienté vers l'intérêt collectif.

Sont strictement interdits :
\begin{itemize}
    \item les mises en scène personnelles ou les contenus dépourvus d'utilité publique ;
    \item les attaques personnelles, provocations, insinuations malveillantes ou éléments relevant de la vie privée d'autrui ;
    \item la republication brute de contenus externes (vidéos, images, extraits, etc.) sans analyse, sans contexte ou sans valeur ajoutée.
\end{itemize}

Ces pratiques n'apportent aucune contribution au débat et vont à l'encontre de la philosophie de la plateforme. Toute violation répétée peut entraîner le retrait du contenu ou des sanctions à l'encontre de son auteur.


\subsection{Les commentaires : prolongement naturel du dialogue}

Les commentaires permettent à la communauté d'interagir avec un Poste, d'enrichir le débat, de préciser un point, de contester une idée ou d'apporter un complément d'information. Ils constituent un prolongement naturel du dialogue et doivent respecter les mêmes exigences que les Postes.

Un commentaire doit être argumenté, respectueux et orienté vers la compréhension ou l'amélioration. Il ne doit jamais contenir d'insultes, de dénigrement, de moqueries ou de qualificatifs dégradants. Les réactions impulsives, les attaques gratuites ou les insinuations malveillantes n'ont pas leur place sur la plateforme. Chaque commentaire doit contribuer à la qualité de l'échange et non à sa dégradation.

\section{Exigences de civilité et de qualité des contributions}

Les règles d'expression définissent le cadre dans lequel les utilisateurs peuvent contribuer à la plateforme. Elles ne restreignent pas la liberté d'expression : elles en garantissent la qualité, la pertinence et le respect mutuel. Leur objectif est de maintenir un environnement sain, constructif et utile pour l'ensemble de la communauté.

\paragraph{Principes généraux.}

Les Postes et les commentaires doivent être rédigés dans un esprit de responsabilité et de respect. Les utilisateurs sont tenus d'éviter tout contenu portant atteinte à la dignité d'autrui, révélant des éléments de vie privée, ou contenant des insultes, propos discriminatoires, attaques personnelles ou tentatives de sabotage. Sont également interdits les contenus mensongers manifestes, les manipulations, les provocations et les polémiques artificielles.

\paragraph{Exigences de qualité des contributions.}

Les contributions doivent être factuelles lorsque nécessaire, argumentées et orientées vers l'intérêt collectif. La plateforme valorise la réflexion, la connaissance et la critique constructive. Chaque intervention doit viser à enrichir la discussion, éclairer un sujet ou contribuer à une compréhension commune.

\paragraph{Sanctions en cas de non-respect.}

Toute violation répétée de ces règles peut entraîner des mesures graduées, allant du retrait du contenu à la suspension temporaire du compte. Ces sanctions visent à protéger l'intégrité de l'espace d'expression et à préserver la qualité des échanges.


Voici une version réorganisée, plus structurée et sans répétitions, tout en conservant l'esprit du texte. Tu peux l'utiliser telle quelle comme section ou sous-section.

\section{Signalement et contestation}

La plateforme met en place un mécanisme de signalement destiné à garantir le respect des règles et la qualité de l'espace d'expression. Le signalement n'a pas vocation à censurer des opinions ni à trancher des désaccords idéologiques : il sert exclusivement à identifier les comportements nuisibles et les violations des principes de bonne conduite.

\paragraph{Mécanisme de signalement.}

Tout utilisateur peut signaler un Poste, un commentaire ou tout autre contenu contrevenant aux règles essentielles de la plateforme. Certains signalements pourront être rendus visibles, notamment depuis la page de l'auteur concerné ou dans une section dédiée aux contenus contestés. Cette transparence vise à prévenir les abus et à permettre à chacun de comprendre les raisons d'un retrait ou d'une sanction.

\paragraph{Procédure de contestation.}

Un utilisateur estimant qu'une décision à son encontre est injustifiée pourra en demander la révision. Une procédure de réexamen sera alors engagée, confiée à un autre groupe de modérateurs chargé d'évaluer la situation de manière indépendante.


\section{Éthique et fonctionnement de la modération}

La modération est assurée par des volontaires issus de la communauté. Leur mission ne consiste pas à juger les idées, opinions ou positions exprimées dans un Poste, ni à déterminer ce qui serait vrai ou faux. Leur rôle est strictement limité à la vérification du respect des règles générales : interdiction des insultes, protection de la vie privée, prévention du sabotage, respect des fondements essentiels du projet et maintien d'une qualité minimale des échanges.

\paragraph{Processus de traitement des signalements.}

Lorsqu'un contenu est signalé, un groupe de modérateurs est tiré au sort parmi les volontaires disponibles. Leur évaluation doit être objective et fondée exclusivement sur les règles établies. Leur responsabilité est de protéger la communauté, non de contrôler la pensée ou d'orienter les débats. Ils doivent agir avec impartialité, transparence et retenue, en évitant tout abus de pouvoir.

\paragraph{Rôle des administrateurs.}

Les administrateurs sont chargés d'attribuer ou de retirer les droits de modération en cas d'abus, de manquement ou de comportement contraire aux valeurs du projet. Ils veillent également au bon fonctionnement général de la plateforme et à la cohérence des décisions prises par les modérateurs.


\section{Conclusion}

Ce chapitre définit le cadre d'expression et d'intégrité des contenus au sein de la plateforme. Le Poste, unité fondamentale d'expression, doit servir l'intérêt collectif et respecter les règles de bonne conduite. Les commentaires doivent suivre la même philosophie. Le signalement et la modération communautaire constituent des mécanismes essentiels pour protéger l'espace, garantir la qualité des échanges et préserver la vision fondatrice du projet. Notre ambition est de créer un environnement où chacun peut s'exprimer librement, mais dans un cadre qui valorise la réflexion, la connaissance et la responsabilité.


% ** A NEW CHAPTER **
\NgelmakProjectNewChapter{Rendre les Déplacements plus Simples pour Tous}[images/transport/jacques-dillies-ITM82DCj6i8-unsplash.jpg][images/transport/joshua-oluwagbemiga-t6nqZ0n3i-k-unsplash.jpg]


\section{Constat et motivations}

La mobilité quotidienne reste l'un des défis les plus pénibles pour une grande partie des populations africaines. Pour aller au travail, rejoindre l'école, assister à un événement ou simplement rendre visite à un proche, il faut souvent traverser un système désorganisé, saturé d'intermédiaires, de pratiques abusives et d'attentes interminables. Le covoiturage, pourtant profondément ancré dans les habitudes depuis des décennies, n'a jamais bénéficié d'une structuration moderne ou d'un accompagnement numérique adapté aux réalités locales. Il fonctionne, mais dans un chaos permanent qui pénalise à la fois les voyageurs et les conducteurs.

Dans de nombreuses villes, la scène est toujours la même : pour trouver un trajet, il faut se rendre physiquement à un point de rassemblement connu de tous. Là, une multitude d'intermédiaires se disputent les voyageurs, les interpellent, les tirent par le bras, crient des destinations et imposent des prix arbitraires. Le conducteur, qui fournit réellement le service, voit sa marge réduite par ces acteurs parasites. Le passager, lui, paie plus cher, attend plus longtemps et subit un stress inutile. Ce système, hérité d'une longue histoire de débrouillardise et d'absence d'organisation publique, est devenu une norme que tout le monde subit sans jamais la remettre en question.

Les tentatives de numérisation existantes, lorsqu'elles apparaissent, se contentent souvent de remplacer les intermédiaires physiques par des plateformes privées qui imposent des commissions exorbitantes. Le résultat est le même : le voyageur paie trop, le conducteur gagne trop peu, et la plateforme s'enrichit. Rien n'est pensé pour les réalités africaines, pour les revenus locaux, pour les contraintes quotidiennes, ni pour la souveraineté numérique. Le problème n'est pas seulement technique : il est structurel, culturel et économique.

\section{Une alternative souveraine et non lucrative}

Face à ce constat, \ngelmak propose une alternative radicalement différente. L'objectif n'est pas de créer une énième plateforme de covoiturage cherchant à maximiser ses profits, mais de construire un service public numérique, souverain, gratuit et orienté vers l'intérêt collectif. Nous refusons la logique entrepreneuriale classique centrée sur le triptyque problème-profit-pour-soi. Notre démarche repose sur une philosophie simple : problème-aide-solution. Nous voulons démontrer qu'un groupe de citoyens peut construire un outil que les gouvernements auraient dû proposer depuis longtemps, un outil qui facilite la vie des gens au lieu de la compliquer, un outil qui aide les jeunes à générer un revenu complémentaire sans être exploités, un outil qui réduit les coûts de transport pour les usagers et fluidifie les déplacements.

Cette vision repose sur un principe fondamental : aucune commission, aucune tarification imposée, aucune exploitation des conducteurs ou des voyageurs. Le service est conçu pour être un bien commun numérique, accessible à tous, et pensé pour répondre aux réalités africaines plutôt qu'aux logiques de profit.

\section{Les usages principaux du service}

\subsection{Le conducteur qui propose un trajet}

Le premier usage du service concerne les conducteurs qui souhaitent proposer un trajet. Un particulier possédant un véhicule ou tout autre moyen de locomotion peut publier un trajet régulier, occasionnel, saisonnier ou spontané. Il peut fixer un prix, laisser la possibilité de négocier ou proposer gratuitement. Les voyageurs intéressés envoient une demande, et le conducteur reste libre d'accepter ou de refuser. Les points de départ et d'arrivée peuvent être fixes ou flexibles, permettant d'ajuster le parcours selon les besoins des passagers.

Ce modèle permet au conducteur de réduire ses coûts de carburant et d'entretien, d'optimiser ses déplacements, de partager la conduite lorsque cela est possible et de générer un revenu complémentaire sans subir la pression d'intermédiaires ou de commissions imposées. Il transforme un trajet déjà prévu en une opportunité d'entraide et de réduction des dépenses.

\subsection{Le voyageur qui initie une demande}

Le second usage concerne les voyageurs qui souhaitent initier une demande de trajet. Un utilisateur peut publier une requête en précisant l'heure, la date, le lieu de départ, le lieu d'arrivée et éventuellement un prix indicatif. Les conducteurs situés dans la zone concernée voient la demande et peuvent l'accepter, proposer un ajustement ou simplement passer leur tour sans aucune pénalité. Le voyageur choisit ensuite l'offre qui lui convient, et les autres propositions disparaissent automatiquement.

Ce modèle permet aux voyageurs de trouver un trajet rapidement, sans se déplacer vers des points de rassemblement, sans subir les cris, les abus ou les attentes interminables. Il redonne au voyageur le contrôle total sur son déplacement et élimine les obstacles inutiles qui rendent les trajets pénibles.

\section{Sécurité, confiance et responsabilité}

La sécurité et la confiance sont au cœur du service. Les conducteurs doivent vérifier leur permis et les documents du véhicule, tandis que les voyageurs peuvent vérifier leur identité. Les données sont cryptées et ne sont utilisées qu'en cas d'incident ou de litige. Un badge indique si un profil est vérifié, et un historique simple permet d'évaluer le sérieux d'un utilisateur. L'objectif n'est pas de créer une surveillance intrusive, mais de garantir un minimum de transparence et de fiabilité pour protéger les usagers.

\section{Une gouvernance communautaire}

La gouvernance du service repose sur une modération communautaire. Des volontaires issus de la communauté traitent les signalements, gèrent les litiges simples et veillent au respect des règles. Ils agissent de bonne foi, en toute neutralité, et transmettent les cas sérieux aux administrateurs. Les administrateurs, quant à eux, assurent la sécurité générale du service, prennent les décisions finales concernant les suspensions ou retraits, et garantissent la transparence du fonctionnement. Ils veillent également à la protection des données et à l'application des politiques définies.

\section{Conclusion}

Ce chapitre présente une vision claire : rendre la mobilité plus simple, plus juste et plus accessible, en s'appuyant sur une solution souveraine, gratuite et pensée pour les réalités africaines. \ngelmak ne cherche pas à exploiter un marché, mais à résoudre un problème. L'objectif est d'offrir un service qui fonctionne, qui aide, qui soulage, et qui démontre que l'Afrique peut construire ses propres outils, adaptés à ses besoins et à ses ambitions. En réorganisant un système déjà existant mais chaotique, en supprimant les intermédiaires abusifs et en redonnant le contrôle aux usagers, nous posons les bases d'une mobilité plus humaine, plus efficace et plus respectueuse de chacun.

\NgelmakProjectNewChapter{Les Prochaines Etapes du Projet \ngelmak}[images/upcoming/nadine-marfurt-kPUQmpH1U-E-unsplash.jpg][images/upcoming/polina-koroleva-YUUSeuO2Hmk-unsplash.jpg]


\section{Perspectives et développements à venir}

\paragraph{Consolider les fondations du Projet \ngelmak.}

Les premiers mois du Projet \ngelmak seront consacrés à la stabilisation des outils essentiels, à l'amélioration de l'expérience utilisateur et à la mise en place d'un environnement technique robuste. Il s'agira d'assurer la fiabilité des services existants, de renforcer la sécurité des données et de garantir une accessibilité optimale, même dans des zones où la connectivité reste limitée. Cette phase de consolidation est indispensable pour bâtir une base solide sur laquelle les futures fonctionnalités pourront s'appuyer.

\paragraph{Étendre progressivement les services.}

Une fois les fondations stabilisées, \ngelmak intégrera de nouveaux services destinés à répondre aux besoins concrets des communautés africaines. L'objectif n'est pas de multiplier les fonctionnalités pour impressionner, mais de proposer des outils réellement utiles, cohérents avec notre vision et adaptés aux réalités du terrain. Les priorités incluront l'amélioration du système de mobilité, l'intégration de services communautaires, la mise en place d'espaces de partage de connaissances et la création de mécanismes permettant aux utilisateurs de contribuer directement à l'évolution de la plateforme.

\paragraph{Renforcer la dimension communautaire.}

\ngelmak n'est pas seulement un outil numérique ; c'est un espace social, un lieu de rassemblement moderne inspiré de la \textit{cour} traditionnelle. Les prochaines étapes viseront à renforcer cette dimension communautaire en facilitant l'organisation de groupes locaux, en encourageant la participation citoyenne et en développant des outils de modération collective. L'objectif est de faire de \ngelmak un espace où chacun peut contribuer, apprendre, transmettre et s'engager dans la construction d'un environnement numérique souverain.

\paragraph{Développer une gouvernance ouverte et transparente.}

À mesure que le projet grandira, il deviendra nécessaire d'établir une gouvernance plus structurée, fondée sur la transparence, la participation et la responsabilité. Les utilisateurs devront pouvoir comprendre comment les décisions sont prises, comment les données sont protégées et comment les orientations futures sont définies. Cette gouvernance ouverte sera un pilier essentiel pour maintenir la confiance, éviter les dérives et garantir que \ngelmak reste un outil au service du bien commun.

\paragraph{Préparer l'interconnexion avec d'autres initiatives africaines.}

Le Projet \ngelmak ne se construit pas en isolation. À long terme, il devra s'interconnecter avec d'autres initiatives africaines, qu'elles soient technologiques, éducatives, culturelles ou économiques. L'objectif est de contribuer à un écosystème numérique africain cohérent, capable de rivaliser avec les grandes plateformes mondiales tout en préservant nos valeurs, nos langues et nos réalités. Cette interconnexion permettra de renforcer la souveraineté numérique du continent et d'ouvrir la voie à de nouvelles formes de collaboration.


\section{Conclusion Générale}

Le Projet \ngelmak s'inscrit dans une démarche de souveraineté, de dignité et de responsabilité collective. Il ne s'agit pas simplement de créer une plateforme numérique, mais de bâtir un espace qui nous appartient, qui reflète nos valeurs et qui répond à nos besoins. Dans un monde où les récits, les technologies et les infrastructures sont souvent contrôlés par d'autres, il devient essentiel de reprendre l'initiative et de construire nos propres outils.

\ngelmak, en tant que \textit{cour} moderne, se veut un lieu de rassemblement, de réflexion et d'action. Un espace où les Africains peuvent se retrouver, partager, apprendre et s'organiser sans dépendre de structures extérieures. Un espace où la mobilité, l'information, la culture et la communauté se rejoignent pour former un réseau souverain, accessible et durable.

Ce document n'est qu'une première étape. Le chemin sera long, parfois difficile, mais nécessaire. Nous avancerons avec humilité, détermination et lucidité. Nous ferons des erreurs, nous corrigerons, nous progresserons. Ce qui compte, c'est la direction : construire un avenir où l'Afrique ne se contente plus de réagir, mais propose, innove et s'impose par ses propres réalisations.

\ngelmak est une invitation à participer à cette construction. Une invitation à reprendre notre place dans le monde numérique. Une invitation à bâtir, ensemble, un espace qui nous ressemble et qui nous rassemble.     

% % ** A NEW CHAPTER **
\NgelmakProjectNewChapter{Livraison et services}
**Développer un Service de Livraison Innovant**


% % ** A NEW CHAPTER **
\NgelmakProjectNewChapter{Lecture, partage et culture}
**Promouvoir la Lecture par l'Échange de Livres**

% % ** A NEW CHAPTER **
\NgelmakProjectNewChapter{Livraison et services}
**Développer un Service de Livraison Innovant**



\end{document}
